\chapter{正文写作方法}

本章展示论文正文中最常用的结构、交叉引用、公式、复杂文字格式、列表和参考文献写法。建议正文只保留内容命令，把版式命令集中放在类文件或少量自定义宏中。

\section{章节层级}

正文使用 \verb|\chapter|、\verb|\section|、\verb|\subsection| 和 \verb|\subsubsection| 组织。模板当前编号为“第1章”、节号为“1.1”，图、表和公式按“章-序号”编号。

\begin{lstlisting}[language=TeX]
\chapter{绪论}
\section{研究背景}
\subsection{国内外研究现状}
\end{lstlisting}

如果需要调整目录深度，可以修改 \texttt{wzu-thesis.cls} 中的设置：

\begin{lstlisting}[language=TeX]
\setcounter{tocdepth}{2}
\end{lstlisting}

如果只是某个小标题不想进入目录，可以使用带星号的标题命令：

\begin{lstlisting}[language=TeX]
\section*{不进目录的标题}
\end{lstlisting}

\section{交叉引用}

模板已经加载 \texttt{cleveref}，推荐使用 \verb|\cref| 自动生成引用名称。写作时先给图、表、公式设置 \verb|\label|，再在正文中引用这些标签。

\begin{lstlisting}[language=TeX]
如 \cref{fig:workflow} 所示，模板把内容与格式分离。
\end{lstlisting}

公式、图和表的标签建议分别使用 \texttt{eq:}、\texttt{fig:} 和 \texttt{tab:} 作为前缀。例如，\cref{eq:sample-loss} 展示了一个常见损失函数。

\section{数学公式}

行内公式可写作 \(a^2+b^2=c^2\)。独立公式建议使用 \texttt{equation} 环境并添加标签：

\begin{equation}
  \mathcal{L}(\theta)=\frac{1}{N}\sum_{i=1}^{N}\left\lVert f_{\theta}(x_i)-y_i\right\rVert_2^2 .
  \label{eq:sample-loss}
\end{equation}

多行推导适合使用 \texttt{align} 环境。等号按 \verb|&| 对齐，每行用 \verb|\\| 换行：

\begin{align}
  \hat{\beta}
  &= \arg\min_{\beta}\left\{(y-X\beta)^{\mathsf T}(y-X\beta)+\lambda\lVert\beta\rVert_2^2\right\} \notag\\
  &= \left(X^{\mathsf T}X+\lambda I_p\right)^{-1}X^{\mathsf T}y .
  \label{eq:ridge}
\end{align}

分段函数可以使用 \texttt{cases}。例如，Huber 损失函数可写为

\begin{equation}
  \rho_{\delta}(r)=
  \begin{cases}
    \dfrac{1}{2}r^2, & |r|\le \delta,\\[0.4em]
    \delta\left(|r|-\dfrac{1}{2}\delta\right), & |r|>\delta .
  \end{cases}
  \label{eq:huber}
\end{equation}

矩阵、向量和转置建议使用统一记号。若正文中采用粗体表示向量和矩阵，可以写作

\begin{equation}
  \mathbf{X}=
  \begin{bmatrix}
    1 & x_{11} & \cdots & x_{1p}\\
    1 & x_{21} & \cdots & x_{2p}\\
    \vdots & \vdots & \ddots & \vdots\\
    1 & x_{n1} & \cdots & x_{np}
  \end{bmatrix},
  \qquad
  \hat{\boldsymbol{\Sigma}}=\frac{1}{n-1}\sum_{i=1}^{n}
  (\mathbf{x}_i-\bar{\mathbf{x}})(\mathbf{x}_i-\bar{\mathbf{x}})^{\mathsf T}.
  \label{eq:matrix}
\end{equation}

复杂优化问题可以把约束条件放在 \texttt{aligned} 子环境中：

\begin{equation}
  \begin{aligned}
    \min_{\mathbf{w},\,b,\,\boldsymbol{\xi}}\quad
    & \frac{1}{2}\lVert\mathbf{w}\rVert_2^2+C\sum_{i=1}^{n}\xi_i\\
    \text{s.t.}\quad
    & y_i(\mathbf{w}^{\mathsf T}\mathbf{x}_i+b)\ge 1-\xi_i,\quad i=1,\ldots,n,\\
    & \xi_i\ge 0,\quad i=1,\ldots,n .
  \end{aligned}
  \label{eq:svm}
\end{equation}

模板已经定义了定理类环境，适合写定义、定理和证明：

\begin{definition}[一致估计]
若估计量序列 \(\hat{\theta}_n\) 满足对任意 \(\varepsilon>0\)，有
\(\lim_{n\to\infty}\Pr(|\hat{\theta}_n-\theta|>\varepsilon)=0\)，则称 \(\hat{\theta}_n\) 是参数 \(\theta\) 的一致估计。
\end{definition}

\begin{theorem}[样本均值的一致性]
设 \(X_1,\ldots,X_n\) 独立同分布，且 \(\mathbb{E}X_i=\mu\)、\(\operatorname{Var}(X_i)=\sigma^2<\infty\)，则样本均值 \(\bar{X}_n\) 是 \(\mu\) 的一致估计。
\end{theorem}

\begin{proof}
由切比雪夫不等式可得
\[
  \Pr(|\bar{X}_n-\mu|>\varepsilon)
  \le \frac{\operatorname{Var}(\bar{X}_n)}{\varepsilon^2}
  = \frac{\sigma^2}{n\varepsilon^2}.
\]
当 \(n\to\infty\) 时，右端趋于 0，因此结论成立。
\end{proof}

正式论文中，公式前后应有必要解释，不建议连续堆叠公式。较长公式如果超出版心，可以优先拆成多行，而不是缩小字号。

\section{复杂文字格式}

正文强调少量关键词时，可以使用 \textbf{黑体强调} 或 \emph{西文斜体强调}。中文论文中不建议大段使用加粗、下划线或彩色文字，除非它们有明确的学术含义。

脚注适合放不影响主线阅读的补充说明\footnote{脚注不宜过长；如果内容很重要，通常应写入正文。}。短引文可直接放在正文中，较长说明可以使用 \texttt{quote} 环境：

\begin{quote}
模板示例中的文字格式只用于说明命令写法。正式论文应保持同一级标题、正文、图题、表题和脚注的格式一致。
\end{quote}

当需要列出概念与说明时，可以使用 \texttt{description}：

\begin{description}
  \item[研究对象] 明确论文分析的主体、范围和时间区间。
  \item[研究方法] 说明模型、算法、数据来源和评价指标。
  \item[研究结论] 对主要发现进行归纳，并避免超出数据支持范围的判断。
\end{description}

代码、命令和文件名可以用等宽字体表示，例如 \texttt{thesis.tex}、\texttt{front/metadata.tex} 和 \verb|latexmk -xelatex thesis.tex|。较长代码建议使用 \texttt{lstlisting}：

\begin{lstlisting}[language=TeX]
\textbf{中文加粗}
\emph{English emphasis}
\footnote{这里写脚注内容。}
\end{lstlisting}

带单位的数字建议使用 \verb|\SI|，例如 \SI{95}{\percent} 置信区间、\SI{30}{\minute} 计算时间和 \SI{1.5}{\giga\byte} 数据文件。这样可以保持数字与单位之间的空距一致。

\section{参考文献}

参考文献条目写在 \texttt{bib/references.bib} 中，正文用 \verb|\cite| 引用。例如，本模板参考了 \LaTeX{} 的基本写作思想\cite{lamport1986latex}，参考文献格式采用 GB/T 7714 风格\cite{gbt7714-2015}。统计模型类论文中常见的文献包括期刊论文\cite{wang2026robust}、会议论文\cite{he2016resnet}、学位论文\cite{li2024thesis}、在线资料\cite{ctan2026biblatex} 和专利\cite{sample-patent}。

模板通过 \texttt{biblatex-gb7714-2015} 自动处理顺序编码、作者名、出版项和标点，日常写作只需要维护 BibTeX 条目。

\begin{lstlisting}[language=TeX]
参考文献可这样引用：\cite{lamport1986latex}
多篇文献可这样引用：\cite{lamport1986latex,gbt7714-2015}
需要把作者放进正文时，可写：Lamport\cite{lamport1986latex} 提出了...
\end{lstlisting}

常见国标参考文献条目写法如下：

\begin{lstlisting}[language=TeX]
@article{demo2026,
  author  = {王一帆 and 赵明},
  title   = {应用统计模型的稳健估计方法},
  journal = {统计研究},
  year    = {2026},
  volume  = {43},
  number  = {6},
  pages   = {1-12}
}

@inproceedings{demo-conf,
  author    = {He, Kaiming and Zhang, Xiangyu},
  title     = {Deep Residual Learning for Image Recognition},
  booktitle = {Proceedings of the IEEE Conference on Computer Vision and Pattern Recognition},
  year      = {2016},
  pages     = {770-778}
}

@mastersthesis{demo-thesis,
  author  = {李明},
  title   = {复杂抽样下居民消费结构统计建模研究},
  school  = {温州大学},
  address = {温州},
  year    = {2024}
}

@online{demo-online,
  author  = {{CTAN Team}},
  title   = {The biblatex Package},
  url     = {https://ctan.org/pkg/biblatex},
  urldate = {2026-06-16}
}
\end{lstlisting}

\section{列表}

有序列表和无序列表使用标准环境即可：

\begin{enumerate}
  \item 先修改 \texttt{front/metadata.tex} 中的论文信息。
  \item 再修改 \texttt{front/abstract.tex} 中的摘要。
  \item 最后在 \texttt{chapters/} 中撰写正文。
\end{enumerate}

无序列表适合放并列说明：

\begin{itemize}
  \item 图题和表题由模板自动居中；
  \item 图形列表和表格列表由 \verb|\makewzulists| 生成；
  \item 参考文献由 \verb|\printbibliography| 输出。
\end{itemize}
